% Appendix: every top-level section starts with \clearpage and uses an apdx:
% label. Section order: references/structure.md (appendix order).
\appendix
\clearpage
\section{Related Works} \label{apdx:related-works}
% ~1 page, 2-3 subsections mapped onto the method's ingredients; each
% paragraph: significance -> taxonomy -> representatives -> shared
% limitation -> bridge. The last subsection ends with the positioning sentence.
\subsection{<Base Paradigm>}
<Field> is now a standard tool for <goal> \citep{Placeholder-A}. Most existing methods fall into two groups, <A> \citep{Placeholder-B} and <B> \citep{Placeholder-C}. Methods of the first group ... Methods of the second group ... Although both groups work well in their settings, they are prone to <problem 1> and <problem 2>.

\subsection{<Enabling Technique>}
<Technique> has recently shown promise for <goal> \citep{Placeholder-A, Placeholder-B}. ... This progress motivates our design of <component>.

\subsection{<Form of Supervision / Constraint>}
... Yet how to <open problem> is still an open question. We close this gap with a framework that <does>, combining <A> with <B>.

\clearpage
\section{Theoretical Analysis} \label{apdx:theory}
% Only when the main text has a formal result [some].
\subsection{Proof of Lemma~\ref{lemma:simplify}} \label{apdx:proof-lemma}
We first rewrite <problem> as <equivalent form> using <identity>. Under <assumption>, <step 1>. Substituting <expression> then gives <step 2>, which is <statement>.

\subsection{Proof of Theorem~\ref{thm:main}} \label{apdx:proof-main-theorem}
\MainTheorem*
\begin{proof}
\textbf{Setup.} Let <symbol> denote <object>.

\textbf{Rewriting the objective.} Under the constraint <condition>, <quantity> becomes:
\begin{align}
    a & = b \notag \\
      & = c,
\end{align}
where the second step follows from <justification>.

\textbf{Bounding the target.} Because <fact>, applying <inequality> gives:
\begin{equation}
    <\text{target quantity}> \;\le\; \mathcal{J}(\theta) + C\,h,
\end{equation}
where $C$ depends only on <constant>. Hence <consequence>.
\end{proof}

\clearpage
\section{Implementation Details} \label{apdx:impl-details}
All experiments use the <framework>\footnote{\url{https://github.com/...}} framework on <hardware>.

\textbf{Training.} <Recipe in prose with the key values: learning rate = <lr>, batch size = <batch size>, ...>. <Why a non-obvious setting was chosen>.

\begin{table}[h]
\centering
\caption{Key hyperparameters used to train X.}
\begin{tabularx}{0.45\linewidth}{l|>{\centering\arraybackslash}X}
\toprule
\multicolumn{1}{c}{\textbf{Hyperparameter}} \vline & \multicolumn{1}{c}{\textbf{Value}} \\
\midrule
Learning rate          & \pend \\
Batch size             & \pend \\
Epochs                 & \pend \\
Max sequence length    & \pend \\
Random seeds           & \pend \\
\bottomrule
\end{tabularx}
\label{tab:hyper-train}
\end{table}

\clearpage
\section{Prompts} \label{apdx:prompts}
\begin{figure}[ht]
\centering
\begin{lrbox}{\PromptBoxA}
\begin{gptpromptbox}
\begin{minipage}{0.95\linewidth}\raggedright\scriptsize\ttfamily
You will rate one training example for difficulty. ... \\
\textcolor{roleB}{<score>} ... \textcolor{roleB}{</score>}
\end{minipage}
\end{gptpromptbox}
\end{lrbox}
\usebox{\PromptBoxA}
\caption{Prompt used for the <role> role.}
\label{fig:prompt-role}
\end{figure}

\clearpage
\section{Additional Experiments} \label{apdx:experiments}
% Subsections mirror the main-text RQs: 1-3 sentences of setup, one table or
% figure, 1-3 sentences of conclusion. Plain headers, generic captions.
\subsection{Results on Other Model Sizes}
We also apply X to <models>. The gains hold at every size we test.

\clearpage
\section{Case Study} \label{apdx:case-study}
\begin{figure}[ht]
\centering
\begin{promptbox}
\small
\textbf{Question:} \quad <question> \\
\textbf{$w/o$ X:} \quad \textcolor{bad}{<failure>} \\
\textbf{$\phantom{/o}w$ X:} \quad \textcolor{good}{<success>}
\end{promptbox}
\caption{A case study of X on <task>.}
\label{fig:case-study}
\end{figure}

% ICLR 2026-style kits only: "Use of LLMs" as the last appendix section.
% ICLR 2027 replaced it with a required AI Use Statement BEFORE the references
% (chapters/8_ethics_reproducibility_statement.tex; references/venues.md §1.5).
% \clearpage
% \section{Use of LLMs} \label{apdx:use-of-llms}
% % 2026-style ICLR kits: always the last appendix section.
% LLMs played no part in forming research ideas, designing experiments, or analysing data. We used them only for:
% \begin{itemize}[leftmargin=*]
%     \item <light editing of a few sentences>;
%     \item <help with plotting scripts>.
% \end{itemize}
% Every result, analysis and conclusion in this paper was produced and checked by the authors, who are fully responsible for its content.
